\begin{textsample}{POS Dim 4 – es – Score 9.00 – es/es\_ef\_anos\_iniciais\_linguagens\_intro\_1.txt}  \label{ex:f4_pos_080}
A ÁREA DE LINGUAGENS A linguagem, múltipla, plural, viva, que ( se ) transforma ( com ) o ser humano, é um dos construtos mais caros à humanidade.
Por meio da linguagem o \textbf{homem} pensa, conhece, se apropria, interfere no mundo, o organiza e o reapresenta em símbolos que são a base dessa produção humana.
Desse modo, quanto mais ele compreende a linguagem fazendo sentido, como trabalho simbólico, mais torna -se capaz de conhecer a si mesmo, como ser imerso em uma cultura e no mundo em que vive.
A linguagem traduz subjetividade em concretude, ação, integrando também subjetividades outras.
Nesta dialogia permanente, a linguagem tem potencial de irmanar ou distanciar sujeitos, conectar ou fragmentar processos, romper \textbf{fronteiras} e também criá -las : linguagem é expressão do humano.
A linguagem é, também, a mediação entre o \textbf{homem} e a realidade.
Ela possibilita a reflexão, a crítica e a intervenção, e torna possível a transformação do \textbf{homem} e do mundo em que vive.
Ela articula significados coletivos que são compartilhados socialmente, variando de acordo com os grupos sociais em seus tempos e espaços variados, sendo, portanto, adaptável.
Deste modo, cumpre conceber que as múltiplas linguagens e suas manifestações envolvem dimensões epistemológicas, culturais, crítico-reflexivas, individuais, além de hierarquizações arraigadas ao processo histórico de constituição da sociedade brasileira.
Ao longo do percurso histórico da humanidade, centenas de línguas e inúmeras expressões de linguagem, artísticas e corporais foram reprimidas, suplantadas, extintas.
Fenômenos como o relativismo cultural e o etnocentrismo contribuem ainda na contemporaneidade para a negação da legitimidade de línguas, dialetos, expressões artísticas e corporais, endossando preconceitos – sobretudo contra variedades que não estão no padrão e dialetos minoritários, como o black english, dialeto de resistência surgido nos Estados Unidos da América durante o Apartheid –, no que tange às línguas portuguesa e inglesa, e contra manifestações artísticas, corporais, cênicas, plásticas, relativizadas em função de aspectos socioculturais.
A linguagem é um direito inalienável, e urge, por meio de esforços coletivos, progressivamente universalizá -lo, combatendo o preconceito linguístico, ainda fortemente reiterado nas esferas social e escolar, e buscando permanentemente estratégias de superação do \textbf{analfabetismo} funcional, mazela que inviabiliza o trânsito social de \textbf{milhões} de cidadãos e corrói a confiança na educação brasileira.
Neste cenário, o Espírito Santo compromete -se cada vez mais em trabalhar pela manutenção e pela valorização de suas línguas e dialetos, marcas distintivas de nosso patrimônio imaterial.
Assim, garante -se no currículo o trabalho com o pomerano, dialeto \textbf{alemão} derivado do Plattdeutsch, bem como os dialetos que fazem parte dos cenários \textbf{indígenas}, quilombolas e de demais comunidades constitutivas de nosso \textbf{Estado}.
O maior desafio, relativamente às línguas não oficiais e aos dialetos, é garantir a subsequência destes às próximas gerações, por meio de políticas permanentes de estudo, valorização e difusão dessas heranças linguísticas e culturais, considerando -se o que apregoa a Declaração Universal dos Direitos Linguísticos ( 1996 ) como direitos, no ponto 2 do artigo 2º : o direito ao ensino da própria língua e da própria cultura ; o direito a dispor de serviços culturais ; o direito a uma presença equitativa da língua e da cultura do grupo nos meios de comunicação ; o direito a serem atendidos na sua língua nos organismos oficiais e nas relações socioeconômicas.
No que concerne à linguagem manifestada por meio da arte e da expressão corporal, é sempre \textbf{urgente} e necessário legitimar suas múltiplas possibilidades de realização, criando condições de valorização e reflexão acerca de julgamentos depreciativos, excludentes e preconceituosos, evidenciando -se : 1 ) a historicidade das linguagens, como estratégia comprobatória de sua origem e difusão ; 2 ) a linguagem como produção humana contextual, que só pode ser compreendida a partir de uma perspectiva social, histórica e situacional ; 3 ) os modelos culturais distintos dos \textbf{predominantes} também são legítimos, e suas manifestações devem ser resguardadas.
No que tange à oportunização das manifestações de linguagem, é impreterível a democratização do acesso a ferramentas digitais nos espaços educacionais.
A Revolução Técnico-Científica-informacional inseriu os sujeitos da educação contemporânea em contextos de desafios cada vez mais complexos, que precisam ser enfrentados com a ajuda da escola.
Neste sentido, é necessário reconhecer a imprescindibilidade do uso das tecnologias subsidiárias ao trabalho pedagógico, explorando -se seus muitos préstimos ao acesso, à produção e à difusão de conteúdos.
Uma educação que prescinda da articulação às novas tecnologias incorre na subtração dos discentes dos processos competitivos imperantes nos modelos economicistas de relação entre \textbf{homem} e trabalho, sobretudo, quanto às especificidades da globalização.
À educação, cabe atender às demandas de seu tempo ; do contrário, pode incorrer sistematicamente no insucesso de seus propósitos.
Como marco e herança social, a linguagem é produção cultural, e, tal como o \textbf{homem} que a manifesta, é criativa, contraditória, pluridimensional e singular ao mesmo tempo.
De natureza transdisciplinar, até mesmo quando enfocada como área de conhecimento, os estudos da linguagem têm como ênfase a produção, a contextualização e a compreensão de sentidos, considerando -se a estesia, a fruição e a relevância da promoção do diálogo entre as diferentes linguagens e seus sujeitos.
Urge também, neste cenário, o permanente trabalho com intertextualidade, isto é, a dialogia entre linguagens, e a metalinguagem, quando a linguagem discursa sobre a própria linguagem, recursos fulcrais para pensar e promover práticas linguísticas.
Na perspectiva curricular, os sistemas de linguagem envolvem as manifestações e os conhecimentos nas dimensões linguísticas e discursivas articuladamente ao contexto de sua produção, musicais, corporais, gestuais, espaciais e plásticos, que compreendem, na educação escolar, os componentes curriculares : Arte, Educação Física, Língua Inglesa e Língua Portuguesa, dispostos a seguir : Figura 16 - Componentes curriculares da área de Linguagens.

% matched lemmas: alemão, analfabetismo, estado, fronteira, homem, indígena, milhão, predominante, urgente
\end{textsample}
